\section{Property Test Catalogue}
\label{sec:appB}
\label{app:properties}

The ten core invariants of Section~\ref{sec:invariants} restated as executable oracles. Each is a precondition that constrains the generated input and a postcondition the operation under test must satisfy. The CDM enumeration universe (Section~\ref{sec:cdm}) supplies the generators: the framework draws random inputs, applies the operation, and checks the postcondition. A failure yields a minimal counterexample. The invariant prose is stated once in Section~\ref{sec:invariants} and not repeated here; this catalogue gives the test forms.

\renewcommand{\arraystretch}{1.3}
\begin{center}
\begin{longtable}{@{}p{2.5cm} p{4.4cm} p{5.4cm}@{}}
\toprule
\textbf{Invariant} & \textbf{Precondition} & \textbf{Postcondition (oracle)} \\
\midrule
\endfirsthead
\toprule
\textbf{Invariant} & \textbf{Precondition} & \textbf{Postcondition (oracle)} \\
\midrule
\endhead

P1 Conservation & $\tau$ applied to state $s$, accepted & $\forall u:\ Q_{s'}(u)=\sum_w w(u)=0$ \\

P2 Atomic commitment & $\tau=\langle m_1,\dots,m_n\rangle$; move $m_k$ fails validation, $1\le k\le n$ & $\texttt{execute}(s,\tau)=s$; no partial state observable \\

P3 Referential integrity & $m$ a move in an accepted $\tau$ & $m.\mathrm{src}\in\mathcal{W}\ \wedge\ m.\mathrm{dst}\in\mathcal{W}\ \wedge\ m.\mathrm{unit}\in\mathcal{U}$ \\

P4 Log monotonicity & log $L=\langle e_1,\dots,e_n\rangle$ after any operation sequence & $\forall i>1:\ e_i.\mathrm{prevHash}=H(e_{i-1})$; no $e_j$, $j<n$, rewritten (append-only, hash-chained) \\

P5 Transaction idempotency & state $s$, transaction $\tau$ & $\texttt{execute}(\texttt{execute}(s,\tau),\tau)=\texttt{execute}(s,\tau)$ (dedup by \texttt{tx\_id}) \\

P6 Lifecycle idempotency & $f(u,\mathit{st},d)=(\mathit{moves},\mathit{st}')$; $\mathit{st}'$ already reflects the event & $f(u,\mathit{st}',d)=(\varnothing,\mathit{st}')$; second call emits no moves \\

P7 Ledger isolation & $m$ a move in virtual ledger $\mathcal{L}_v$ & $m.\mathrm{src}\in\mathcal{L}_v\ \wedge\ m.\mathrm{dst}\in\mathcal{L}_v$; no move crosses the virtual/real boundary \\

P8 Snapshot consistency & time $t$ in the log; $s_t$ the state at $t$ & $\forall w,u:\ \texttt{clone\_at}(t).\texttt{balance}(w,u)=s_t.\texttt{balance}(w,u)\ \wedge\ \texttt{clone\_at}(t).\texttt{unit\_state}(u)=s_t.\texttt{unit\_state}(u)$ \\

P9 Lifecycle purity & lifecycle function $f$, fixed inputs $(u,\mathit{st},d)$ & repeated evaluation yields identical $(\mathit{moves},\mathit{st}')$; no effect beyond the return \\

P10 PnL path-independence & price vectors $P_{t_0},P_{t_1}$; any transaction sequence on $[t_0,t_1]$ & $\mathrm{PnL}_{\text{price}}+\mathrm{PnL}_{\text{flow}}=V_{t_1}-V_{t_0}$, independent of the intermediate path \\

\bottomrule
\end{longtable}
\end{center}
\renewcommand{\arraystretch}{1.0}

\begin{remark}[P8 and the projection discipline]
The P8 oracle reconstructs \emph{both} balances and unit state from the log. \texttt{unit\_state} is read from UnitStatus, a materialised projection of the immutable event log---a read cache the log always rebuilds, never a separate source of truth (Section~\ref{sec:lifecycle}). The oracle is the executable form of that discipline: a catamorphism---replaying the events up to $t$ rebuilds the exact value that held then---so $\texttt{clone\_at}(t).\texttt{unit\_state}(u)=s_t.\texttt{unit\_state}(u)$ holds by construction, and a failure exposes any drift between the cache and the log.
\end{remark}

\noindent\textbf{Totality oracles.} Two global oracles the same generator exercises, not numbered among the core ten:
\begin{itemize}
\item \emph{Totality.} For every generated input $i$, $\texttt{execute}(s,i)$ returns success or a well-typed error---never an unhandled exception.
\item \emph{Valid transitions only.} For every $(\text{product type},\text{state},\text{event})$ triple absent from the product's transition table, the lifecycle function returns an error. This is the executable witness that illegal lifecycle transitions are unrepresentable (Section~\ref{sec:lifecycle}).
\end{itemize}

\noindent\textbf{PnL inputs.} $V_t$ and the price vectors $P_t(u)$ that drive P10 are produced by the valuation projection; their construction and the state-aware price function are specified in the valuation/market-data spec, not duplicated here.

\noindent\textbf{SBL and obligation extensions.} Section~\ref{sec:sbl-invariants} adds ten invariants P11--P20 specific to Securities Borrowing and Lending (collateralisation, lender ownership invariance, partial-return and settlement monotonicity), and Section~\ref{sec:obligation-invariants} adds P21--P23 for obligation liveness. These never relax a core invariant: every transaction satisfying P11--P23 also satisfies P1--P10. Their executable forms follow the same precondition/postcondition shape.

\medskip
\noindent\textbf{Typed oracles.} Each oracle above has a typed counterpart: a precondition that constrains the generated input and a postcondition the operation under test must satisfy. Both are total, so the generator and the checker share one type. The operation under test returns a typed \texttt{Outcome}; \texttt{Accepted} is the \emph{only} constructor that carries a ledger, so a partially applied (``half-committed'') ledger is not representable. That absence is the \emph{Totality} oracle and the ``no partial state'' half of P2 made structural rather than tested.

\begin{lstlisting}[language=Haskell]
-- execute = the canonical (Sec. 4) three-home transaction-apply, lifted to
-- return Outcome; supplied by the implementation, not a symbol defined here.
data Outcome = Accepted Ledger | Rejected TxError   -- no third "crashed" case

-- A property: a precondition constraining the generated input and a
-- postcondition over the outcome of running the operation on it. Both total.
data Property i o = Property
  { pre  :: Ledger -> i -> Bool
  , post :: Ledger -> i -> o -> Bool }

data TxInput = TxInput { tiTx :: OracleTx, tiScope :: Scope }
data Scope   = Scope   { scWallets :: [WalletId], scUnits :: [UnitId] }

registered :: Ledger -> UnitId -> Bool
registered l u = maybe False (const True) (productTerms l u)

-- Sum a wallet's holding of u across a wallet set. Summation is a monoid
-- homomorphism into the Qty group -- the empty set sums to mempty (=0) with no
-- special case, so a zero-holder unit conserves vacuously (no divide-by-holders).
unitTotal :: Ledger -> [WalletId] -> UnitId -> Qty
unitTotal l ws u = foldMap (\w -> maybe mempty psBalance (position l w u)) ws

-- Equality of ledgers BY OBSERVATION over the scope -- the only equality the
-- oracles need; the abstract Ledger exports no structural Eq.
sameLedger :: Scope -> Ledger -> Ledger -> Bool
sameLedger sc a b =
     and [ position a w u == position b w u | w <- scWallets sc, u <- scUnits sc ]
  && and [ unitStatus a u == unitStatus b u | u <- scUnits sc ]
\end{lstlisting}

\noindent Five oracles test the same operation, \texttt{execute}, and share the one \texttt{Property TxInput Outcome} shape. Each postcondition that asserts a positive fact is guarded on \texttt{Accepted}, so it claims nothing about a rejected transaction.

\begin{lstlisting}[language=Haskell]
-- P1 Conservation. On any accepted transaction every unit's holdings sum to
-- zero across the scope (closed system); a rejected one asserts nothing.
p1 :: Property TxInput Outcome
p1 = Property (\_ _ -> True) $ \_ (TxInput _ sc) o -> case o of
  Rejected _ -> True
  Accepted l -> all (\u -> unitTotal l (scWallets sc) u == mempty) (scUnits sc)

-- P2 Atomic commitment. Any move failing validation rejects the whole txn;
-- no partial ledger exists (Outcome cannot carry one).
p2 :: Property TxInput Outcome
p2 = Property (\l (TxInput tx _) -> any (invalidMove l) (otMoves tx))
              (\_ _ o -> isRejected o)

-- P3 Referential integrity. Every move in an accepted txn names a registered
-- unit and in-scope wallets.
p3 :: Property TxInput Outcome
p3 = Property (\_ _ -> True) $ \_ (TxInput tx sc) o -> case o of
  Rejected _ -> True
  Accepted l -> all (\m -> registered l (mUnit m)
                        && inScope sc (mFrom m) && inScope sc (mTo m)) (otMoves tx)

-- P7 Ledger isolation. With Scope = the virtual ledger's wallets, no move
-- crosses its boundary -- a property of the moves, independent of the outcome.
p7 :: Property TxInput Outcome
p7 = Property (\_ (TxInput tx _) -> not (null (otMoves tx)))
     $ \_ (TxInput tx sc) _ ->
         all (\m -> inScope sc (mFrom m) && inScope sc (mTo m)) (otMoves tx)

-- P5 Transaction idempotency (dedup by tx_id). Re-running an accepted txn is a
-- no-op; this oracle re-runs the operation, so it takes it directly.
idempotentTx :: (Ledger -> OracleTx -> Outcome)
             -> Scope -> Ledger -> OracleTx -> Bool
idempotentTx exec sc l tx = case exec l tx of
  Rejected _  -> True
  Accepted l' -> case exec l' tx of Accepted l'' -> sameLedger sc l' l''
                                    Rejected _   -> False
\end{lstlisting}

\noindent The remaining five test different operations -- the log append, the snapshot rebuild, the pure lifecycle handler, the PnL decomposition. Forcing them into the \texttt{Property} shape would buy nothing, so each is a bespoke total predicate (the restraint rule: stop at the cheaper guarantee).

\begin{lstlisting}[language=Haskell]
-- P4 Log monotonicity. The log is append-only and hash-chained BY CONSTRUCTION
-- (built only by `append`, which stamps the predecessor's hash; no rewrite or
-- delete is exported). The oracle re-checks the link; a break needs API-tamper.
data Hashed = Hashed { hPrev :: Hash, hPayload :: Transaction, hSelf :: Hash }
chainOK :: [Hashed] -> Bool
chainOK es = and (zipWith (\a b -> hPrev b == hSelf a) es (drop 1 es))

-- P8 Snapshot consistency. cloneAt t rebuilds the state at t by replaying the
-- events up to t (a catamorphism -- replaying the log always rebuilds the exact
-- value that held then), so snapshot must equal s_t by observation (both
-- balances and unit state -- sameLedger covers both).
p8 :: Scope -> Ledger -> Ledger -> Bool          -- cloneAt t  vs  s_t
p8 = sameLedger

-- A lifecycle handler is a PURE TOTAL function. Purity (no IO in the type) IS
-- P9: re-evaluation cannot differ and there is no effect to leak.
type Lifecycle st ev = UnitId -> st -> ev -> ([Move], st)

-- P6 Lifecycle idempotency. Feeding the handler its own output state emits no
-- further moves and fixes the state.
idempotentL :: Eq st => Lifecycle st ev -> UnitId -> st -> ev -> Bool
idempotentL f u st ev = let (_, st') = f u st ev in f u st' ev == ([], st')

-- Valid-transitions-only. For every (product, state, event) absent from the
-- transition table, the guarded handler returns an error: not in table => Left.
-- An illegal transition is thus not representable as success.
validTransitionsOnly
  :: (p -> st -> ev -> Bool)                         -- in the table?
  -> (p -> st -> ev -> Either TxError ([Move], st))  -- the guarded handler
  -> p -> st -> ev -> Bool
validTransitionsOnly inTable h p st ev
  | inTable p st ev = True
  | otherwise       = case h p st ev of Left _ -> True; Right _ -> False

-- P10 PnL path-independence. V_t and the price vectors P_t(u) come from the
-- valuation projection (valuation/market-data spec); the decomposition closes
-- regardless of the intermediate transaction path.
p10 :: Cash -> Cash -> Cash -> Cash -> Bool      -- pnlPrice pnlFlow V_t0 V_t1
p10 pnlPrice pnlFlow v0 v1 = pnlPrice <> pnlFlow == v1 <> cashNeg v0
\end{lstlisting}

\noindent The full self-contained listing (helpers \texttt{inScope}, \texttt{isRejected}, \texttt{invalidMove}, the \texttt{Move}/\texttt{OracleTx} types, and the \texttt{total} coverage witness) reuses the FORMALIS-cleared \texttt{Ledger} library and its total accessors, consolidated in \texttt{reference/Ledger.hs}.
